\documentclass[a4paper,12pt]{article}
\usepackage[T2A]{fontenc}
\usepackage[utf8]{inputenc}
\usepackage[russian]{babel}
\usepackage{amsmath,amssymb,mathtools}
\usepackage{helvet}
\renewcommand{\sfdefault}{cmss}
\renewcommand{\familydefault}{\sfdefault}
\usepackage{icomma}
\usepackage[margin=23mm,top=28mm,bottom=24mm]{geometry}
\usepackage{microtype}
\usepackage{tikz}
\usetikzlibrary{arrows.meta,positioning,shapes.geometric,calc}
\usepackage{pgfplots}
\pgfplotsset{compat=1.18}
\usepackage{tabularx,booktabs}
\usepackage{enumitem}
\usepackage{hyperref}
\usepackage{parskip}
\usepackage{fancyhdr}
\usepackage{setspace}
\usepackage{xcolor}
\usepackage{titlesec}
\definecolor{accent}{HTML}{26777A}
\definecolor{darkaccent}{HTML}{174A51}
\definecolor{textgray}{HTML}{343C40}
\definecolor{muted}{HTML}{687477}
\definecolor{lightline}{HTML}{D8E3E2}
\color{textgray}
\hypersetup{hidelinks,pdftitle={Чек-лист. Повторение профильного ЕГЭ},pdfauthor={Лёвин Артём Александрович}}
\setstretch{1.19}
\setlength{\parindent}{0pt}
\setlength{\parskip}{7pt}
\titlespacing*{\section}{0pt}{19pt}{8pt}
\titlespacing*{\subsection}{0pt}{13pt}{5pt}
\titleformat{\section}{\normalfont\Large\bfseries\color{darkaccent}}{\thesection.}{.6em}{}
\titleformat{\subsection}{\normalfont\large\bfseries\color{accent}}{\thesubsection.}{.6em}{}
\setlength{\headheight}{25pt}
\pagestyle{fancy}
\fancyhf{}
\fancyhead[L]{\small\color{muted} Математика · ЕГЭ, профильный уровень}
\fancyhead[R]{\small\color{muted} 09.10.2026}
\fancyfoot[C]{\small\color{muted}\thepage}
\renewcommand{\headrulewidth}{0.3pt}
\renewcommand{\headrule}{\hbox to\headwidth{\color{lightline}\leaders\hrule height \headrulewidth\hfill}}
\newcommand{\key}[1]{\textbf{\color{darkaccent}#1}}
\newcommand{\fine}[1]{{\small\color{muted}#1}}
\newcommand{\answerline}{\par\vspace{8mm}\noindent\color{lightline}\rule{0.78\linewidth}{0.35pt}\color{textgray}\par}
\begin{document}
\pagenumbering{Alph}
\begin{titlepage}
  \thispagestyle{empty}
  \begin{center}
    \vspace*{26mm}
    {\color{accent}\Large МАТЕМАТИКА · ПРОФИЛЬНЫЙ ЕГЭ\par}
    \vspace{19mm}
    {\fontsize{31}{39}\selectfont\bfseries\color{darkaccent}Чек-лист\par}
    \vspace{10mm}
    {\fontsize{21}{27}\selectfont Повторение: алгебра,\par
      прикладные и текстовые задачи\par}
    \vspace{20mm}
    {\large 9 октября 2026 года\par}
    \vspace{6mm}
    {\normalsize Теория · разобранные примеры · 10 заданий · ответы\par}
    \vfill
    {\normalsize\color{darkaccent}Лёвин Артём Александрович\\[3pt]
      эксклюзивно для Ксении Клыковой\par}
    \vspace{28mm}
  \end{center}
  \begin{tikzpicture}[remember picture,overlay]
    \draw[accent!55,line width=1.2pt] ([xshift=-6mm,yshift=23mm]current page.south west)
      .. controls ([xshift=45mm,yshift=38mm]current page.south west)
        and ([xshift=-35mm,yshift=15mm]current page.south east)
      .. ([xshift=6mm,yshift=26mm]current page.south east);
  \end{tikzpicture}
\end{titlepage}
\pagenumbering{arabic}

\section{Алгебраическая разминка}
\subsection{Линейное уравнение}
\key{Правило.} В уравнении $a-b=c$ неизвестное вычитаемое находится как $b=a-c$. Следите за знаками при переносе слагаемых.
\[
28-2x=4\quad\Longrightarrow\quad 2x=28-4=24\quad\Longrightarrow\quad\boxed{x=12}.
\]
В уравнениях с дробями сначала укажите значения, при которых знаменатели обращаются в нуль. Умножайте обе части на общий знаменатель только с учётом этих ограничений.

\subsection{Степень суммы и действия со степенями}
\key{Правила:} $a^m a^n=a^{m+n}$; $\dfrac{a^m}{a^n}=a^{m-n}$ при $a\ne0$; $(a^m)^n=a^{mn}$.
\[
(x+8)^5=243=3^5\ \Longrightarrow\ x+8=3\ \Longrightarrow\ \boxed{x=-5}.
\]
Поскольку пятая степень нечётная, равенство пятых степеней действительных чисел влечёт равенство оснований. \key{Нельзя} распределять степень по сумме: $(x+8)^5\ne x^5+8^5$ в общем случае.

\key{Пример на сокращение степеней} (при $x\ne0$):
\[
\frac{2(3x^7)^2x^6}{9x^{20}}
=\frac{2\cdot9x^{14}\cdot x^6}{9x^{20}}=\boxed{2}.
\]
Каждый множитель в скобках возводится в степень; при умножении одноимённых оснований показатели складываются.

\subsection{Подстановка и общий множитель}
Если $a=3b$, найдите значение дроби
\[
\frac{a+9b+16}{a+3b+8}
=\frac{12b+16}{6b+8}
=\frac{4(3b+4)}{2(3b+4)}=\boxed{2}.
\]
\key{Ограничение:} исходный знаменатель $6b+8\ne0$, поэтому $b\ne-\frac43$. Сокращать общий множитель можно лишь тогда, когда он отличен от нуля.

\clearpage
\section{Прикладные формулы: определить величину и единицы}
\subsection{Расстояние до горизонта}
При $L=\sqrt{2Rh}$, где $R$ и $h$ выражены в километрах, из $L=16$~км и $R=6400$~км получаем
\[
16^2=2\cdot6400\cdot h,\qquad h=\frac{256}{12800}=0{,}02\ \text{км}=\boxed{20\ \text{м}}.
\]
\key{Проверка:} $h\ge0$; сантиметры, метры и километры нельзя смешивать в одной формуле.

\subsection{Определить, что дано, а что требуется}
Если по условию расстояние $x=80$, а искомая величина задаётся формулой
\[
y=0{,}023x^2-0{,}5x+32,
\]
то \(y=0{,}023\cdot6400-40+32=\boxed{139{,}2}\). Прочитайте обозначения перед подстановкой: число $80$ относится к $x$, а не к $y$.

\section{Два корня: сначала смысл, потом ответ}
\subsection{Когда мяч выше заданной отметки?}
Высота мяча (в метрах) в момент $t$ (в секундах) задана законом
\[
h(t)=2+14t-5t^2.
\]
Найдём промежуток, когда мяч находится \key{выше 10 м}:
\[
2+14t-5t^2>10\ \Longleftrightarrow\ 5t^2-14t+8<0.
\]
Граничные моменты — корни уравнения $5t^2-14t+8=0$:
\[
D=196-160=36,\qquad t_1=0{,}8,\qquad t_2=2.
\]
Поскольку ветви параболы $h(t)$ направлены вниз, условие $h(t)>10$ выполнено при
\[
\boxed{0{,}8<t<2},\qquad\text{продолжительность }\boxed{2-0{,}8=1{,}2\ \text{с}}.
\]
Момент $0{,}8$~с соответствует подъёму через отметку 10 м, а $2$~с — опусканию через неё. Граничные моменты \key{не входят} в интервал строгого неравенства.

\begin{center}
\begin{tikzpicture}
\begin{axis}[
 width=.9\linewidth,height=72mm,
 xmin=0,xmax=2.8,ymin=0,ymax=14,
 axis lines=left,xlabel={$t$, с},ylabel={$h$, м},
 ylabel style={at={(axis description cs:-.07,.96)},anchor=north},
 xtick={0,0.8,1.4,2,2.8},xticklabels={$0$,$0{,}8$,$1{,}4$,$2$,$2{,}8$},
 ytick={0,2,10,12},grid=major,grid style={lightline},
 tick label style={font=\small},label style={font=\small},
 enlargelimits=false]
 \addplot[domain=0:2.8,samples=180,very thick,color=accent]{2+14*x-5*x^2};
 \addplot[domain=0:2.8,dashed,color=muted]{10};
 \addplot[only marks,mark=*,mark size=2.2pt,color=darkaccent] coordinates {(0.8,10)(2,10)};
 \node[anchor=south,font=\small,color=darkaccent] at (axis cs:1.38,11.8){выше 10 м};
\end{axis}
\end{tikzpicture}
\end{center}
\fine{На рисунке показана математическая модель высоты; для реального движения используется лишь физически допустимое время.}

\subsection{Отбор корней в контексте}
\begin{itemize}[leftmargin=*,itemsep=2pt,topsep=3pt]
 \item \key{Время и скорость:} проверяйте неотрицательность времени, положительность скорости и ограничения знаменателей.
 \item \key{Два момента:} различайте вход в заданную область и выход из неё; если спрашивают длительность, находите разность моментов.
 \item \key{До и после изменения уровня:} при свободном падении из покоя $s=\frac{gt^2}{2}$; подъём воды равен разности прежней и новой высоты падения. Например, при $g\approx10\ \text{м/с}^2$, $t_{\text{до}}=2$~с, $t_{\text{после}}=1{,}6$~с: $\Delta h=5(2^2-1{,}6^2)=7{,}2$~м.
\end{itemize}

\clearpage
\section*{Алгоритм: отбор корней в прикладной задаче}
\thispagestyle{fancy}
\vspace{4mm}
\begin{center}
\begin{tikzpicture}[
 >=Stealth, node distance=10mm,
 start/.style={ellipse,draw=accent,very thick,align=center,minimum width=42mm,minimum height=12mm,font=\small},
 blocknode/.style={rectangle,rounded corners=2pt,draw=lightline!60!darkaccent,fill=white,
   align=center,text width=115mm,minimum height=15mm,font=\small,inner sep=5pt},
 decision/.style={diamond,aspect=3.9,draw=accent,very thick,align=center,
   text width=62mm,inner sep=3pt,font=\small},
 side/.style={rectangle,rounded corners=2pt,draw=lightline!60!darkaccent,
   align=center,text width=28mm,minimum height=14mm,font=\small,inner sep=3pt},
 flow/.style={-{Stealth[length=2.4mm]},line width=.85pt,color=accent}]
\node[start] (start) {Начало};
\node[blocknode,below=of start] (read) {Определить искомую величину и её единицы.\\Записать ограничения: ОДЗ, $t\ge0$, $v>0$ и т.~п.};
\node[blocknode,below=of read] (solve) {Составить уравнение и найти все математические корни.};
\node[decision,below=15mm of solve] (check) {Подходит ли корень\\по условиям задачи?};
\node[side,right=11mm of check] (discard) {Отбросить этот корень};
\node[blocknode,below=17mm of check] (keep) {Сохранить допустимые корни.\\Проверить таким же способом остальные.};
\node[blocknode,below=of keep] (meaning) {Ответить именно на вопрос: момент времени,\\длительность, скорость, масса или другая величина.};
\node[start,below=of meaning] (finish) {Ответ};
\draw[flow] (start)--(read);
\draw[flow] (read)--(solve);
\draw[flow] (solve)--(check);
\draw[flow] (check)--node[left,font=\small,color=darkaccent]{да}(keep);
\draw[flow] (check)--node[above,font=\small,color=darkaccent]{нет}(discard);
\draw[flow] (discard.south) |- ([yshift=4mm]keep.east) -- (keep.east);
\draw[flow] (keep)--(meaning);
\draw[flow] (meaning)--(finish);
\end{tikzpicture}
\end{center}
\clearpage

\section{Текстовые задачи: одна логика, разные сюжеты}
\subsection{Движение по реке}
Пусть собственная скорость лодки $v$~км/ч, скорость течения $1$~км/ч. При движении по $255$~км в каждую сторону против течения затрачено на $2$~ч больше, чем по течению.
\[
\frac{255}{v-1}-\frac{255}{v+1}=2,\qquad v>1.
\]
\[
\frac{510}{v^2-1}=2\ \Longrightarrow\ v^2=256\ \Longrightarrow\ v=\pm16.
\]
\key{Отбор:} собственная скорость $v>1$, поэтому $\boxed{v=16\ \text{км/ч}}$. Значение $-16$ не имеет физического смысла. Корень $v=1$ тоже был бы недопустим: знаменатель обращается в нуль.

\subsection{Сплавы и концентрации}
Первый сплав содержит $5\%$ меди, второй — $13\%$; второй тяжелее первого на $9$~кг. Массовая доля меди в объединённом сплаве $10\%$.
\begin{center}
\renewcommand{\arraystretch}{1.28}
\begin{tabularx}{\linewidth}{@{}Xccc@{}}
\toprule
Сплав & Масса, кг & Доля меди & Масса меди, кг \\
\midrule
Первый & $x$ & $0{,}05$ & $0{,}05x$\\
Второй & $x+9$ & $0{,}13$ & $0{,}13(x+9)$\\
Общий & $2x+9$ & $0{,}10$ & $0{,}10(2x+9)$\\
\bottomrule
\end{tabularx}
\end{center}
\[
0{,}05x+0{,}13(x+9)=0{,}10(2x+9).
\]
Получаем $0{,}18x+1{,}17=0{,}20x+0{,}90$, откуда $x=13{,}5$~кг. Значит, масса общего сплава $2\cdot13{,}5+9=\boxed{36\ \text{кг}}$.

\subsection{Совместная работа}
Два рабочих вместе выполняют заказ за $2$ дня. За один день первый делает столько, сколько второй за два дня. Значит, \key{производительность первого вдвое больше}, чем второго: $p_1=2p_2$.
\[
p_1+p_2=\frac12\quad\Longrightarrow\quad 3p_2=\frac12,
\quad p_2=\frac16,\quad p_1=\frac13.
\]
Если полный заказ принят за $1$, первый выполнит его самостоятельно за $1/p_1=\boxed{3\ \text{дня}}$.
\key{Важно:} часть работы, выполненная за определённое время, задаёт отношение производительностей. Её абсолютный объём для этого вывода знать не требуется.

\clearpage
\section*{Тренировочный блок · Путь А}
\fine{10 заданий. Решайте последовательно. В прикладных задачах указывайте единицы измерения и обосновывайте отбор корней.}
\begin{enumerate}[label=\textbf{\arabic*.},leftmargin=*,itemsep=5mm,topsep=6mm]
\item Решите уравнение: $35-3x=8$.
\answerline
\item Решите уравнение: $(x+4)^5=32$.
\answerline
\item Упростите выражение при $x\ne0$:
  $\displaystyle \frac{(2x^4)^2\cdot3x^5}{6x^{13}}$.
\answerline
\item Найдите значение дроби $\displaystyle\frac{a+4b+12}{a+b+6}$,
  если $a=2b$ и исходная дробь определена.
\answerline
\item По формуле $L=\sqrt{2Rh}$ вычислите высоту $h$ в метрах,
  если $R=5000$~км и $L=20$~км.
\answerline
\newpage
\item Высота мяча (в метрах) задаётся формулой $h(t)=2+12t-4t^2$,
  где $t$ --- время в секундах. Сколько секунд мяч находится на высоте
  \emph{строго выше} $10$~м? Укажите интервал времени.
\answerline
\item Камень отпускают без начальной скорости с одной и той же высоты
  над водой. До подъёма уровня воды он падал $2$~с, после подъёма ---
  $1{,}6$~с. Найдите, на сколько метров поднялась вода, считая
  $s=5t^2$ (метры, секунды).
\answerline
\item Лодка прошла $24$~км против течения и $24$~км по течению.
  Скорость течения $1$~км/ч; путь против течения занял на $2$~ч больше.
  Найдите собственную скорость лодки. Объясните отбор корней.
\answerline
\item Первый сплав содержит $6\%$ меди, второй --- $14\%$.
  Масса второго на $5$~кг больше. После соединения массовая доля
  меди стала $11\%$. Найдите массу полученного сплава.
\answerline
\item Два рабочих вместе выполняют заказ за $3$ дня. Первый за день
  выполняет вдвое больше работы, чем второй. За сколько дней
  первый выполнит весь заказ самостоятельно?
\answerline
\end{enumerate}
\vspace{5mm}
\textbf{Проверьте перед сдачей:} определены ли переменные, указаны ли ОДЗ\newline
и единицы измерения, объяснён ли отбор корней и дан ли ответ на вопрос.
\clearpage
\section*{Ответы к тренировочному блоку}
\renewcommand{\arraystretch}{1.56}
\begin{center}
\begin{tabularx}{\linewidth}{@{}cX@{}}
\toprule
№ & Ответ и контрольное условие \\
\midrule
1 & $x=9$. \\
2 & $x=-2$. \\
3 & $2$ при $x\ne0$. \\
4 & $2$ при $b\ne-2$. \\
5 & $40$ м ($0{,}04$ км). \\
6 & $1<t<2$; длительность $1$ с. \\
7 & $7{,}2$ м. \\
8 & $5$ км/ч; корень $-5$ отвергается, поскольку $v>1$. \\
9 & $20$ кг. \\
10 & $4{,}5$ дня. \\
\bottomrule
\end{tabularx}
\end{center}
\vspace{6mm}
\key{Самопроверка:} в задачах 6 и 8 двух корней недостаточно для ответа — необходимо установить их физический смысл; в задаче 7 вычисляется разность расстояний.
\end{document}
